\section{Charge-preserving preparation of the asymptotic end}
\label{en:section}

This section constructs the strict rest exteriors used in the comparison
argument. The construction preserves the two closed flux forms, controls
the area of every full cut, and relates the resulting energy to the
invariant mass of the original end. All constraint densities in this
section are geometric densities, as in Definition~\ref{def:data}:
\[
 2\mu=R_g+\tau^2-|K|_g^2,
 \qquad J=\operatorname{div}_g(K-\tau g).
\]
Thus the factors in the ADM fluxes remain $16\pi$ and $8\pi$.

\subsection{Flux forms and conformal changes}

Set $X_1=\mathcal E$, $X_2=\mathcal B$, and
$\alpha_i=\iota_{X_i}dV_g$. The divergence constraints say precisely
that $d\alpha_i=0$. For any other oriented metric $h$ on $\Omega$,
the same forms determine fields by $\iota_{X_i(h)}dV_h=\alpha_i$.
Their two fluxes through every full cut remain $4\pi Q_E$ and
$4\pi Q_B$. This statement uses Stokes' Theorem in the truncated
exterior of the cut, including each portion coincident with $S$.

For $|Y|_g\leq1$ define
\begin{equation}\label{en:ball-tests}
 \begin{split}
 I_g(Y)&=|X_1|_g^2+|X_2|_g^2
             +2\langle X_1\times X_2,Y\rangle_g,\\
 D_g(Y)&=\mu+J(Y)-I_g(Y)
         =\mu_{\mathrm m}+J_{\mathrm m}(Y).
 \end{split}
\end{equation}
The inequality $2|X_1\times X_2|\leq |X_1|^2+|X_2|^2$ gives
$I_g(Y)\geq0$, and
\[
 \inf_{|Y|_g\leq1}D_g(Y)=\mu_{\mathrm m}-|J_{\mathrm m}|_g.
\]
Consequently charged DEC implies total DEC. A constant $SO(2)$ rotation
of $(\alpha_1,\alpha_2)$ preserves both $I_g$ and $D_g$: it preserves
the sum of the squared norms and the cross product.

\begin{lemma}[Conformal identities]\label{en:conformal}
Let $u>0$, $g'=u^4g$, $K'=u^2K$, and retain $\alpha_1,\alpha_2$.
Then
\begin{equation}\label{en:conformal-constraints}
 \begin{split}
 u^4\mu'&=\mu-4u^{-1}\Delta_g u,\\
 u^2J'&=J+4K(\nabla\log u,\cdot),\\
 X_i(g')&=u^{-6}X_i(g),\qquad
 I_{g'}(u^{-2}Y)=u^{-8}I_g(Y).
 \end{split}
\end{equation}
In particular,
\begin{equation}\label{en:conformal-exact}
 \begin{split}
 u^4D_{g'}(u^{-2}Y)-D_g(Y)
 ={}&4u^{-1}\{-\Delta_g u+K(\nabla u,Y)\}\\
    &+(1-u^{-4})I_g(Y).
 \end{split}
\end{equation}
If $u\geq1$, it follows that
\begin{equation}\label{en:conformal-margin}
 u^4D_{g'}(u^{-2}Y)
 \geq D_g(Y)+4u^{-1}(-\Delta_g u-|K|_g|du|_g).
\end{equation}
For a fixed oriented surface and either fixed sign $\sigma\in\{-1,1\}$,
write $\theta_\sigma=H+\sigma\operatorname{tr}_S K$. Then
\begin{equation}\label{en:conformal-expansion}
 \theta'_\sigma=u^{-2}(\theta_\sigma+4\partial_\nu\log u).
\end{equation}
\end{lemma}

\begin{proof}
The scalar-curvature formula in dimension three is
$R_{u^4g}=u^{-5}(-8\Delta_g u+R_gu)$. Both $\tau^2$ and
$|K|^2$ acquire the factor $u^{-4}$. This proves the first identity.
Write $\pi=K-\tau g$ and $f=\log u$. The transformed trace reversal
is $\pi'=u^2\pi$. In the divergence of this tensor use
\[
 (\nabla'_A-\nabla_A)B
 =2\{df(A)B+df(B)A-g(A,B)\nabla f\}.
\]
Contracting the two covariant tensor slots gives
$u^2\operatorname{div}_{g'}(u^2\pi)
=\operatorname{div}_g\pi+4K(\nabla f,\cdot)$; the terms
containing $\tau\,df$ cancel. This proves the momentum identity.
Since $dV_{g'}=u^6dV_g$, fixed flux forms give $X_i'=u^{-6}X_i$.
In an orthonormal frame their components therefore scale by $u^{-4}$,
whereas $u^{-2}Y$ has the same orthonormal components as $Y$.
The electromagnetic identity and Equation~\ref{en:conformal-exact}
follow. Its last term is nonnegative for $u\geq1$, and
$K(\nabla u,Y)\geq-|K||du|$ proves
Equation~\ref{en:conformal-margin}. Finally
$H'=u^{-2}(H+4\partial_\nu\log u)$ and
$\operatorname{tr}'_S K'=u^{-2}\operatorname{tr}_S K$.
\end{proof}

\subsection{ADM flux under a change of asymptotic plane}

\begin{lemma}[Asymptotic Lorentz flux]\label{en:adm-flux}
Let $G=\eta+h$ be a Lorentzian metric on an exterior region containing
the spacelike planes under consideration and the portions of their
connecting cylinders at large radius. In background inertial
coordinates suppose, uniformly on these regions,
\[
 h=O_2(r^{-q_0}),\qquad q_0>1/2,\qquad
 \operatorname{Ein}(G)=O(r^{-4}).
\]
Use signature $(-+++)$ and the positive second-form convention
$K(U,V)=G(\nabla_U n,V)$. If the ADM limits on one plane exist,
they exist on every fixed uniformly spacelike plane, and they are the
components of the same translation covector $\mathcal C$:
\begin{equation}\label{en:translation-covector}
 \mathcal C(a)=E a^0+\sum_iP_i a^i.
\end{equation}
In particular, if $E>|P|$ and $m=(E^2-|P|^2)^{1/2}$, the unit
future vector $b=(E/m,-P/m)$ satisfies $\mathcal C(b)=m$.

The same conclusion applies to a stationary vacuum or electrovacuum
end whenever these estimates hold in its stationary coordinates,
after a constant linear change making its limiting metric Minkowski.
In coordinates adapted to a stationary field with unit timelike
limit, suppose the shift has order $O_1(r^{-q_0})$.
Then the energy of a rest plane equals that of the spatial orbit
metric. This energy is also unchanged by spatial coordinate changes
$y=x+\zeta(x)$ with
$\zeta=O_2(r^{1-q_1})$, $q_1>1/2$.
\end{lemma}

\begin{proof}
Indices in the following calculation are raised with $\eta$.
Define
\begin{equation}\label{en:superpotential}
 \begin{split}
 \mathcal F_{lab}={}&\partial_a h_{lb}-\partial_l h_{ab}
   +\eta_{lb}(\partial^d h_{ad}-\partial_a\operatorname{tr}_\eta h)\\
   &-\eta_{ab}(\partial^d h_{ld}-\partial_l\operatorname{tr}_\eta h).
 \end{split}
\end{equation}
It is skew in $l,a$. Expanding its divergence and collecting the
two Hessians, wave operator, and trace terms of the linearized Ricci
tensor gives
\begin{equation}\label{en:superpotential-divergence}
 \partial^l\mathcal F_{lab}=2\operatorname{Ein}^{\mathrm{lin}}_{ab}(h)
   =O(r^{-2-2q_0})+O(r^{-4}).
\end{equation}
The last estimate follows because the difference between exact and
linearized Einstein curvature is bounded by
$C(|h||\partial^2h|+|\partial h|^2)$.

In coordinates adapted to a background plane $x^0=0$, direct
substitution in Equation~\ref{en:superpotential} gives
\begin{equation}\label{en:superpotential-components}
 \begin{split}
 \mathcal F_{i00}&=\partial_jh_{ij}-\partial_ih_{jj},\\
 \mathcal F_{i0k}&=2(K^{\mathrm{lin}}_{ik}
                -\delta_{ik}\operatorname{tr}K^{\mathrm{lin}})
                +\partial_kh_{0i}-\delta_{ik}\partial_jh_{0j},\\
 K^{\mathrm{lin}}_{ik}
   &=\tfrac12(\partial_0h_{ik}-\partial_ih_{0k}
                                      -\partial_kh_{0i}).
 \end{split}
\end{equation}
Here $j$ is spatial. The last expression for $K^{\mathrm{lin}}$
follows by expanding the future unit normal and its connection;
in Gaussian normal coordinates it reduces to
$K^{\mathrm{lin}}=\tfrac12\partial_0h$.
For fixed $k$, the extra vector in the middle line has identically
zero spatial divergence. Its flux through any exterior coordinate
sphere is zero: extend $h_{0i}$ smoothly across the ball and integrate
that identity. The value of the extension cannot affect the flux.
The nonlinear errors in $K$ and in its trace reversal are
$O(r^{-1-2q_0})$ and hence have integrated flux
$O(R^{1-2q_0})$. Consequently the plane fluxes
$\int\mathcal F_{i0b}n^i\,dA/(16\pi)$ equal $(E,P_i)$.

For a fixed translation vector $a$, regard
$\tfrac12\mathcal F_{lab}a^b dx^l\wedge dx^a$ as a two-form and
take its background Lorentzian Hodge dual. With orientation
$dx^0\wedge dx^1\wedge dx^2\wedge dx^3$, its integral on an
outward plane sphere is the flux just computed, with the common
orientation chosen to give positive Schwarzschild energy. Equation
\ref{en:superpotential-divergence} bounds its exterior derivative.
Two large plane sections can be joined along a cylinder in
$r\asymp R$ with absolute three-dimensional coordinate volume
$O(R^3)$. Stokes' Theorem bounds the difference of the integrals by
\[
 C(R^{1-2q_0}+R^{-1})\longrightarrow0.
\]
Replacing the sections by spheres or uniformly comparable ellipsoids
in their planes has the same estimate, using the intervening annular
three-region. The limiting flux is thus independent of the plane for
fixed $a$. Tensor covariance under constant Lorentz changes proves
Equation~\ref{en:translation-covector}, including the plus sign of its
spatial components. Substitution of $b=(E/m,-P/m)$ gives $m$.

For the last assertions, write the stationary metric in a rest
coordinate system as
\[
 G=\gamma_{ij}dy^i dy^j+2\beta_i dy^i dt-Wdt^2,
 \qquad W\longrightarrow1,\quad\beta=O_1(r^{-q_0}).
\]
The orbit metric is $h_s=\gamma+W^{-1}\beta\otimes\beta$, so
$h_s-\gamma=O_1(r^{-2q_0})$. Its extra ADM flux is
$O(R^{1-2q_0})$ and vanishes. Finally let $\gamma=\delta+k$ with
$k=O_2(r^{-q_0})$, and express its pullback under $y=x+\zeta(x)$.
Taylor expansion with one derivative gives
\[
 \gamma_{\mathrm{new}}
 =\delta+k+2\operatorname{sym}D\zeta+\mathcal R,
 \qquad
 \mathcal R=O_1(r^{-2q_1}+r^{-q_0-q_1}).
\]
The linear ADM integrand of $2\operatorname{sym}D\zeta$ is
$\Delta\zeta_i-\partial_i\operatorname{div}\zeta
=\partial_j(\partial_j\zeta_i-\partial_i\zeta_j)$.
Its flux is zero by a smooth interior extension and skew symmetry.
The remainder flux is
$O(R^{1-2q_1}+R^{1-q_0-q_1})\to0$. This proves coordinate
invariance in the precise class used here. A change of stationary
time coordinate changes the threading one-form but leaves $h_s$
itself unchanged.
\end{proof}

\subsection{A strict conformal direction and the timelike endpoint}

\begin{lemma}[Strict conformal direction]\label{en:strictification}
Suppose $g-\delta=O_2(r^{-q})$ and $K=O_1(r^{-1-q})$, $q>1/2$,
on the single end. Fix $0<\delta_0<\min(q,1)$. There are smooth
functions $b\geq |K|_g$, $a>0$, $w>0$, and $\phi>0$ such that
\begin{equation}\label{en:strict-pde}
 \begin{gathered}
 -\Delta_g\phi=b\sqrt{|d\phi|_g^2+a^2}+w,
 \qquad \partial_\nu\phi=-1\quad\hbox{on }S,\\
 \phi=O_2(r^{-1}),\qquad
 b=Cr^{-1-q},\quad a=r^{-2},\quad w=r^{-3-\delta_0}
 \quad\hbox{far out}.
 \end{gathered}
\end{equation}
The gradient flux of $\phi$ has a finite limit at infinity. In
particular $u=1+s\phi$, $s>0$, preserves charged DEC and makes it
strict, with margin at least $c_s r^{-3-\delta_0}$ on the tail,
and makes every weak fixed-sign boundary expansion strictly negative.

If $K$ is compactly supported and $g-\delta=O_\ell(r^{-1})$ for
every $\ell$, one may choose $b$ compactly supported. In that case
$-\Delta_g\phi=w$ on the tail and $\phi=O_\ell(r^{-1})$ for every
$\ell$. If also $R_g=O_\ell(r^{-4})$, then
$R_{(1+s\phi)^4g}=O_\ell(r^{-3-\delta_0})$ for every $\ell$.
\end{lemma}

\begin{proof}
Choose smooth majorants and positive interpolations with the stated
tails; the bound on $K$ permits a constant radial majorant $Cr^{-1-q}$
outside a compact set. When $K$ is compactly supported, choose $b$ with
compact support instead. On $\Omega_L=\{r\leq L\}$, including the
compact core, impose the additional boundary condition $\phi=0$ at
$r=L$. Continue from $t=0$ to $t=1$ in
\[
 -\Delta_g\phi=t b\sqrt{|d\phi|_g^2+a^2}+w,
 \qquad \partial_\nu\phi=-1,\quad \phi|_{r=L}=0.
\]
The linearization has principal part $-\Delta_g$ and drift
$-t b\nabla\phi/\sqrt{|d\phi|^2+a^2}$ of norm at most $b$.
The inner Neumann and outer Dirichlet boundaries are disjoint smooth
components. The homogeneous mixed problem has zero kernel: a
nonconstant positive maximum or negative minimum is excluded by the
strong maximum principle and the boundary point lemma, and a constant
is excluded by the Dirichlet condition. The mixed elliptic Fredholm
index is zero, by continuation from the mixed Laplacian, so the
linearization is invertible.

Here are bounds closing this continuation. On a fixed $\Omega_L$,
mixed $W^{2,p}$ estimates, $p>3$, and
$\sqrt{|d\phi|^2+a^2}\leq |d\phi|+a$ give
\[
 \|\phi\|_{W^{2,p}}
 \leq C_L(1+\|d\phi\|_{L^p}+\|\phi\|_{L^p})
 \leq \tfrac12\|\phi\|_{W^{2,p}}
          +C'_L(1+\|\phi\|_{L^p}).
\]
The second inequality is gradient interpolation, with its small
parameter chosen after $C_L$. If the supremum norms were unbounded,
divide a sequence by its supremum norm. Compactness in $C^1$ gives
a function of supremum one solving
$-\Delta z=t b|dz|$, with zero mixed boundary data. This is a
homogeneous equation with bounded drift, by taking the drift to be
$t b\nabla z/|dz|$ where $dz\ne0$ and zero elsewhere. The same
maximum principles give a contradiction. The resulting $W^{2,p}$
bounds, followed by elliptic regularity, close the continuation.
The solution is nonnegative: an interior negative minimum contradicts
$-\Delta\phi>0$, and a negative minimum at $S$ contradicts the boundary
point lemma, since the outward domain derivative is
$-\partial_\nu\phi=1$.

The bounds can be made independent of $L$ on each compact subset.
To see this without a bounded-domain constant depending on $L$, put
$f(r)=r^{-1}(1-r^{-\delta_0})$ on a sufficiently distant tail.
Direct differentiation gives
\begin{equation}\label{en:tail-barrier}
 -\Delta_g f-|K|_g|df|_g
 =\delta_0(1+\delta_0)r^{-3-\delta_0}
      +O(r^{-3-q}).
\end{equation}
The same estimate holds with $b$ in place of $|K|$. Since
$ba=O(r^{-3-q})$ and $\delta_0<q$, a sufficiently large multiple
$A f$ satisfies
$-\Delta_g(Af)\geq b\sqrt{|d(Af)|^2+a^2}+w$.
Comparison on the exterior annulus, choosing its inner boundary
value to dominate $\phi$, gives
\begin{equation}\label{en:exhaustion-tail}
 0\leq\phi\leq C(1+\sup_{\Omega_{L_*}}\phi)r^{-1}
 \quad(r\geq L_*),
\end{equation}
with fixed $L_*$. If the core supremum diverged along an exhaustion,
normalize by it and apply the preceding local estimates, including
the inner boundary estimates. A subsequence converges locally in
$C^1$ to a nonzero nonnegative solution of
$-\Delta z=b|dz|$, with $\partial_\nu z=0$ and $z=O(r^{-1})$.
The tail estimate forces its positive maximum to occur in a compact
set; the strong and boundary maximum principles again exclude it.
Thus exhaustion gives a positive smooth solution of
Equation~\ref{en:strict-pde}.

Rescale each exterior annulus to unit size. The metric has uniformly
bounded $C^2$ coefficients, the scaled drift has norm $O(r^{-q})$,
and the already proved bound is $\phi=O(r^{-1})$. Interior
$W^{2,p}$ estimates and then Schauder estimates give
$d\phi=O(r^{-2})$, $d^2\phi=O(r^{-3})$.
In the compact-$b$ case the equation is the linear Poisson equation
on a smooth symbol end; differentiating its rescaled form gives the
estimates at every order. Its right side is integrable in either
case. The divergence theorem therefore gives a finite limit of
$\int_{r=L}\partial_{\nu_L}\phi\,dA_g$. Its difference from the
Euclidean gradient flux is $O(L^{-q})$, so that limit also exists.
For $u=1+s\phi$,
$-\Delta u-|K||du|\geq s w$. Apply
Equations~\ref{en:conformal-margin} and
\ref{en:conformal-expansion}. The scalar-curvature conclusion follows
from its conformal formula. The energy increment is finite and
tends to zero with $s$, because its ADM flux contribution is
$-(s/2\pi)\lim\int\partial_{\nu_L}\phi\,dA_g$.
\end{proof}

\begin{lemma}[Reduction of the non-timelike endpoint]\label{en:timelike-reduction}
Suppose the numerical conclusion of Theorem~\ref{thm:numerical} has
been proved for all data in its class satisfying $E>|P|$.
Then every datum in the full class satisfies $E>|P|$, and the
numerical conclusion holds for the full class.
\end{lemma}

\begin{proof}
Choose $0<\delta_0<\min(q,1)$ and $f$ as in
Equation~\ref{en:tail-barrier}. Far enough out $f'<0$ and
$-\Delta_g f-|K||df|\geq0$. Let $\chi(r)$ be a smooth nondecreasing
cutoff from zero to one supported in this valid range. Define
\[
 \phi(r)=\int_r^\infty\chi(t)(-f'(t))\,dt
\]
there, and extend it constantly over the compact core. It is bounded,
positive, smooth, and equals $f$ far out. The radial differentiation
formula shows
\[
 -\Delta_g\phi-|K||d\phi|
 =\chi(-\Delta_g f-|K||df|)-\chi' f'|dr|_g^2\geq0.
\]
For each $s\geq0$ set $u_s=1+s\phi$ and use
$(u_s^4g,u_s^2K)$ with the original forms. Charged DEC and all
boundary signs persist. Completeness persists, since $u_s\geq1$,
and the full-cut infimum can only increase. The data still have
decay exponent greater than $1/2$. The new constraint terms have
orders $r^{-3-\delta_0}$ and $r^{-3-q}$, so their absolute
integrability follows from that of the old sources and
Equation~\ref{en:conformal-constraints}.

The asymptotic expansion
$u_s=1+s/r+o_2(r^{-1})$ gives, directly in the original chart,
\begin{equation}\label{en:energy-raising}
 E_s=E+2s,\qquad P_s=P,\qquad Q_{E,s}=Q_E,
 \qquad Q_{B,s}=Q_B.
\end{equation}
Indeed the leading new metric term is $4s\delta_{ij}/r$ and its
energy flux is $32\pi s$; the factor $1/(16\pi)$ yields $2s$.
The trace reversal is multiplied by $u_s^2$, whose additional
momentum flux is $O(r^{-q})$ and tends to zero.

The assumed timelike numerical conclusion implies
$r_{A,s}\leq2\sqrt{(E+2s)^2-|P|^2}$. By
Proposition~\ref{ct:full-cut-positive}, $A_{\min,g}(S)>0$, while
$r_{A,s}\geq\sqrt{A_{\min,g}(S)/(4\pi)}$. If $E\leq|P|$,
take $s\downarrow(|P|-E)/2$ from above. The right side tends to
zero and contradicts this fixed positive lower bound. Thus
$E>|P|$. Applying the assumed timelike result to the original
data finishes the proof. No common limiting boost is used.
\end{proof}

\subsection{Compact corrections for the linearized constraints}

On Euclidean three-space write
\[
 (\mathcal LH)=\partial_i\partial_j
                (H_{ij}-\operatorname{tr}H\,\delta_{ij}),
 \qquad
 (\mathcal Dk)_i=\partial_j
                (k_{ij}-\operatorname{tr}k\,\delta_{ij}).
\]
The formal adjoint kernels are respectively the affine functions
and the Euclidean Killing fields. For the scalar operator this follows
from $\partial_i\partial_jp-(\Delta p)\delta_{ij}=0$: taking the
trace gives $\Delta p=0$, then its Hessian vanishes. For momentum,
the adjoint equation is equivalent to the Euclidean Killing equation
after taking its trace. Thus the momentum tests are translations
and rotations, with no dilation test.

\begin{lemma}[Supported linear inverses]\label{en:compact-inverses}
Fix a compact subannulus of $\{1<|z|<4\}$. A smooth scalar source
$f$ supported there has a compactly supported symmetric solution
$H$ of $\mathcal LH=f$ if its affine moments vanish. A smooth
vector source $V$ has a compactly supported symmetric solution
$k$ of $\mathcal Dk=V$ if its Killing moments vanish. The solutions
can be supported in a fixed larger compact subannulus, and, for
each integer $\ell\geq0$ and $1<p<\infty$,
\begin{equation}\label{en:inverse-estimates}
 \|H\|_{W^{\ell+2,p}}\leq C\|f\|_{W^{\ell,p}},
 \qquad
 \|k\|_{W^{\ell+1,p}}\leq C\|V\|_{W^{\ell,p}}.
\end{equation}
\end{lemma}

\begin{proof}
We first specify the ordinary divergence inverse used in the proof.
For a mean-zero smooth compactly supported source in a ball, the
support-preserving Bogovski\u\i\ operator gives a smooth compactly
supported vector field, with divergence equal to the source and a
gain of one derivative in $W^{\ell,p}$. The regularized integral
operators and their support and Sobolev mapping properties are
given in \cite[Theorems~3.2 and~4.9, Corollary~4.11]{CostabelMcIntosh2010}.
To use balls on the annulus,
cover the prescribed source support by finitely many balls whose
closures lie inside the annulus and join them by a finite connected
family of such balls. Split a source by a partition of unity.
Along a spanning tree of overlapping balls transfer the integral
of each leaf piece to its parent using a fixed unit-integral bump
in the overlap. Removing leaves successively leaves each piece
with integral zero, since the total integral is zero. Applying the
ball operators and adding the results gives a divergence inverse
with fixed compact support. The finite cutoffs and bumps preserve
all the stated Sobolev estimates. Use successively larger fixed
compact supports for successive applications below.

For a scalar $f$ with zero affine moments solve $\partial_iV_i=f$.
Compact support and integration by parts imply
$\int V_i=-\int z_i f=0$. Solve
$\partial_j B_{ij}=V_i$ componentwise. Then the symmetric part
$P=(B+B^{\mathsf T})/2$ satisfies
$\partial_i\partial_jP_{ij}=f$. In dimension three the inverse
of trace reversal is
\[
 H=P-\tfrac12\operatorname{tr}P\,\delta,
 \qquad H-\operatorname{tr}H\,\delta=P.
\]
This gives the scalar solution and gains two derivatives.

For momentum, zero translation moments permit a componentwise solve
$\partial_jB_{ij}=V_i$. Put $S_{ij}=(B_{ij}-B_{ji})/2$.
The rotation moments give
\[
 0=\int(z_iV_j-z_jV_i)=2\int S_{ij}.
\]
Solve $\partial_kD_{ijk}=-S_{ij}$, choosing $D$ skew in $i,j$,
and define
\[
 C_{ijk}=D_{ijk}-D_{ikj}-D_{jki}.
\]
Directly using that skew symmetry gives
$C_{ijk}=-C_{ikj}$ and
$(C_{ijk}-C_{jik})/2=D_{ijk}$. It follows that
$P_{ij}=B_{ij}+\partial_kC_{ijk}$ is symmetric and
$\partial_jP_{ij}=V_i$. Set
$k=P-\tfrac12\operatorname{tr}P\,\delta$.
There is one net derivative of gain: the second divergence solve
gains a derivative and the displayed correction differentiates
it once. All supports remain compact in the prescribed larger
subannulus. Integration against the adjoint kernels also proves
the necessity of the stated moment conditions.
\end{proof}

\begin{lemma}[Weak-end annular replacement]\label{en:annular-replacement}
Let $(g,K)$ have the weak end bounds of Definition~\ref{def:data},
integrable geometric constraint densities, and total DEC. Let
$(g_*,K_*)$ be smooth vacuum reference data on that end with
$g_*-\delta=O_\ell(r^{-1})$, $K_*=O_\ell(r^{-2})$ at every order,
and the same ADM energy and momentum. For arbitrarily large $R$
there are smooth data $(\widehat g_R,\widehat K_R)$ equal to the
original data on $r\leq R$ and to the reference on $r\geq4R$,
such that on the intermediate annulus
\begin{equation}\label{en:annular-deficit}
 \widehat\mu_R-|\widehat J_R|_{\widehat g_R}
 \geq-\eta_RR^{-3},\qquad \eta_R\longrightarrow0.
\end{equation}
For every fixed $1/2<\beta<\min(q,1)$ the scaled perturbations
from Euclidean data are $O(R^{-\beta})$ in $C^2\times C^1$.
\end{lemma}

\begin{proof}
Use $z=x/R$ on $A=\{1<|z|<4\}$ and regard the scaled fields as
$g(Rz)$ and $RK(Rz)$. The constraints scale by $R^2$. Write
$H=g(Rz)-g_*(Rz)$ and $k=R(K(Rz)-K_*(Rz))$. These have
$C^2\times C^1$ norm $O(R^{-\beta})$. If $\lambda$ equals one
near the inner boundary, zero near the outer boundary, and takes
values in $[0,1]$, blend the metric and second form using $\lambda$.
The linearized constraint error beyond the blend of the two linear
sources is
\[
 f_R=[\mathcal L,\lambda]H,\qquad V_R=[\mathcal D,\lambda]k.
\]
If $P(H)=H-\operatorname{tr}H\,\delta$, the scalar commutator is
\begin{equation}\label{en:commutator-orders}
 f_R=(\partial_i\partial_j\lambda)P(H)_{ij}
           +2(\partial_i\lambda)\partial_jP(H)_{ij},
 \qquad
 (V_R)_i=(\partial_j\lambda)P(k)_{ij}.
\end{equation}
Thus $\|f_R\|_{W^{1,p}}+\|V_R\|_{W^{1,p}}=O(R^{-\beta})$
for every fixed finite $p$. Only the second derivatives of the
original metric and first derivatives of its second form are used.

We next estimate all adjoint-kernel moments. Put
$h=g-g_*$, $\kappa=K-K_*$ in the physical $x$ coordinates, and
$s=\mathcal L_xh$, $V=\mathcal D_x\kappa$. They belong to $L^1$.
In fact the scalar linear source for either datum equals its
$2\mu$ plus an error bounded by
\[
 C\bigl(|g-\delta||\partial^2g|+|\partial g|^2+|K|^2\bigr),
\]
and the momentum linear source equals $J$ plus an error bounded by
\[
 C\bigl(|g-\delta||\partial K|+|\partial g||K|\bigr).
\]
These errors are $O(r^{-2-2q})$ for the original data and
$O(r^{-4})$ for the reference. They are integrable because $q>1/2$.
For any integrable source $F$ on the end,
\begin{equation}\label{en:sublinear-moment}
 \int_{r\leq4R}r|F|\,dx=o(R).
\end{equation}
Indeed splitting at a fixed $M$ and dividing by $R$ bounds the
outer contribution by $4\int_{r>M}|F|$, while the inner contribution
tends to zero; then let $M\to\infty$.

For an affine scalar $p$ define its linear constraint boundary flux
\[
 B_p(t;h)=\int_{r=t}
  \{p\,\partial_jP(h)_{ij}-(\partial_jp)P(h)_{ij}\}
          n^i\,dA.
\]
For a Euclidean Killing field $Z$, define
$B_Z(t;\kappa)=\int_{r=t}P(\kappa)_{ij}Z_i n^j\,dA$.
Integration by parts gives, for either type of test, the difference
of the two boundary fluxes as the integral of the test times the
linear source; the adjoint-kernel term is zero. For $p=1$ and
constant $Z$, the boundary flux differences tend to zero as
$t\to\infty$, since the two ADM vectors agree. For momentum the
difference between Euclidean and geometric trace reversal has
integrated error $O(t^{1-2q})+O(t^{-1})\to0$.
For homogeneous linear $p$ or rotational $Z$, integration from a
fixed sphere and Equation~\ref{en:sublinear-moment} instead give
\begin{equation}\label{en:first-moment-flux}
 B_p(t;h)=o(t),\qquad B_Z(t;\kappa)=o(t).
\end{equation}
No limit for these first-moment fluxes is required.

For clarity, the scaling factors in the commutator moments are as
follows. A constant test gives
\[
 \int_A f_R\,dz
 =-R^{-1}\left\{B_1(R;h)
       +\int_{R<r<4R}\lambda(x/R)s(x)\,dx\right\}=o(R^{-1}).
\]
A homogeneous linear test $p(z)$ gives
\[
 \int_Ap(z)f_R(z)\,dz
 =-R^{-2}\left\{B_{p(x)}(R;h)
       +\int_{R<r<4R}\lambda(x/R)p(x)s(x)\,dx\right\}
 =o(R^{-1}).
\]
For momentum the identical calculation uses $R\kappa(Rz)$;
translation tests have factor $R^{-1}$, rotation tests factor
$R^{-2}$. Equations~\ref{en:sublinear-moment} and
\ref{en:first-moment-flux} therefore give
\begin{equation}\label{en:small-cokernel}
 \int_Ap f_R=o(R^{-1}),\qquad
 \int_A Z\cdot V_R=o(R^{-1})
\end{equation}
for every fixed affine $p$ and every fixed Killing field $Z$.

Choose a fixed nonnegative bump $\chi$ supported in the annulus
and positive on a ball. The Gram matrix of the affine basis with
weight $\chi$ is positive definite. Subtract the linear combination
of the smooth functions $\chi p$ having the same affine moments
as $f_R$. Likewise use the Gram matrix of the Killing-field basis
and the vector bumps $\chi Z$ for $V_R$. Equation
\ref{en:small-cokernel} says that the coefficients of all these
bumps are $o(R^{-1})$. Their norms of every fixed order therefore
have this size. The remaining sources have zero moments and
$W^{1,p}$ norm $O(R^{-\beta})$. Apply
Lemma~\ref{en:compact-inverses} with the negative signs to cancel
them. The resulting compact corrections have norms
\[
 \|H_R^{\mathrm c}\|_{W^{3,p}}
       +\|k_R^{\mathrm c}\|_{W^{2,p}}=O(R^{-\beta}).
\]
Taking $p>3$ controls their $C^2\times C^1$ norms and ensures
positivity of the corrected metric for large $R$.

The nonlinear constraint map about $(\delta,0)$ has quadratic
remainder bounded in $C^0$ by $C R^{-2\beta}$ on these scaled
data. Its sources consequently differ from the convex blend of
the original and vacuum sources by
$o(R^{-1})+O(R^{-2\beta})$ in $C^0$. The original scaled
momentum is $O(R^{-\beta})$; changing its norm from the original
metric to the corrected metric costs only $O(R^{-2\beta})$.
Convexity of the cone $\{(\mu,J):\mu\geq|J|\}$ therefore
proves a scaled DEC deficit of size
$o(R^{-1})+O(R^{-2\beta})=o(R^{-1})$.
Undoing the scaling multiplies sources by $R^{-2}$ and gives
Equation~\ref{en:annular-deficit}. All corrections are supported
away from the two annular boundaries, so the asserted exact
matching and smoothness follow.
\end{proof}

\subsection{The strict rest exterior}

\begin{proposition}[Rest preparation]\label{en:rest-preparation}
Let the data satisfy Definition~\ref{def:data} and $E>|P|$.
Assume $\theta_+(S)\leq0$ on every boundary component. More generally,
one may prescribe a fixed sign $\sigma\in\{-1,1\}$ on each component
and assume $H+\sigma\operatorname{tr}_S K\leq0$ there.
Set $m=\sqrt{E^2-|P|^2}$. There are smooth data
$(g_j,K_j,\mathcal E_j,\mathcal B_j)$ on the same oriented exterior
$\Omega$, with the same closed forms $\alpha_1,\alpha_2$, having
the following properties.
\begin{enumerate}
 \item Both charges are exactly $Q_E,Q_B$. Charged DEC is strict
 everywhere; for some $c_j>0$ and a fixed $0<\delta_0<1$,
 \[
  \inf_{|Y|_{g_j}\leq1}D_{g_j}(Y)
       \geq c_j r_y^{-3-\delta_0}
       \quad\hbox{on the far end}.
 \]
 Every boundary component has strictly negative future expansion,
 or strictly negative prescribed expansion
 $H+\sigma\operatorname{tr}_S K$ in the fixed-sign variant.
 \item In a rest chart $y$, $g_j-\delta=O_\ell(r_y^{-1})$ and
 $R_{g_j}=O_\ell(r_y^{-3-\delta_0})$ for every $\ell$.
 The tensor $K_j$ is compactly supported, and the two flux forms
 have components $O_1(r_y^{-2})$. The constraint densities are
 integrable and the metric is complete with $S$ included.
 \item The rest ADM momentum is zero and its energy satisfies
 $E_j\to m$. The fields and tensors converge to the original data
 smoothly on every compact subset of $\Omega$.
 \item There are $\epsilon_j\downarrow0$ such that
 $g_j\geq(1-\epsilon_j)g$ on all of $\Omega$, and
 \begin{equation}\label{en:area-convergence}
  A_{\min,g_j}(S)\longrightarrow A_{\min,g}(S).
 \end{equation}
\end{enumerate}
No uniform bound on the compact geometry of this sequence is asserted
or needed. Each prepared datum is fixed before the subsequent
elliptic comparison parameters are chosen.
\end{proposition}

\begin{proof}
\emph{The reference plane.}
Since $m>0$, use the exterior Schwarzschild metric
\begin{equation}\label{en:schwarzschild-reference}
 G_*=-\left(\frac{1-m/(2r_y)}{1+m/(2r_y)}\right)^2dT^2
                  +(1+m/(2r_y))^4|dy|^2
\end{equation}
at sufficiently large radius. Set $v=P/E$, so $|v|<1$, and
$A=(I-vv^{\mathsf T})^{-1/2}$. Parametrize the plane
$T=v\cdot y$ by $y=Ax$. Its induced background Minkowski metric
is exactly $|dx|^2$. The unit normal and an adapted spatial frame
are related to the horizontal ones by the Lorentz boost with
velocity $v$. By Lemma~\ref{en:adm-flux}, the plane data have
\[
 (E_*,P_*)=(\gamma m,\gamma m v)=(E,P),
 \qquad\gamma=(1-|v|^2)^{-1/2}.
\]
They are vacuum, with $g_*-\delta=O_\ell(r_x^{-1})$ and
$K_*=O_\ell(r_x^{-2})$ at every order. Identify their $x$
coordinates with the original end coordinates and apply
Lemma~\ref{en:annular-replacement}.

\emph{Bending the reference end.}
Beyond $r_x=4R$ replace the plane by the graph
\begin{equation}\label{en:bend-height}
 T(y)=(v\cdot y)\psi\left(\frac{\log(r_y/R_b)}{L}\right),
\end{equation}
where $\psi$ is smooth, is one on $(-\infty,0]$, is zero on
$[1,\infty)$, and has values in $[0,1]$. Choose $L$ sufficiently
large that
$|v|(1+\|\psi'\|_\infty/L)<1$; then choose $R_b$ to be a
sufficiently large fixed multiple of $R$. The bend starts in the
exact reference portion. Its background gradient has norm at
most the preceding strict bound, so it is uniformly spacelike in
$G_*$ once $R$ is large. It becomes a horizontal Schwarzschild
slice for $r_y\geq R_b e^L$. Denote the resulting data before
conformal repair by $(g_R,K_R)$.

For fixed $L$, the graph has $|dT|\leq c<1$,
$|d^2T|\leq C/r_y$, and derivatives of the background Schwarzschild
coefficients are $O(r_y^{-2})$. The graph formulas thus give
uniform comparability of $g_R$ with $\delta_y$,
$|\partial g_R|\leq C/r_y$, and $|K_R|\leq C/r_y$ throughout
the bend. The annular corrections have stronger corresponding
bounds. It follows that, throughout these large-radius ranges,
\begin{equation}\label{en:radial-bounds}
 c\leq |dr_y|_{g_R}\leq C,
 \qquad |\Delta_{g_R}r_y|\leq C/r_y,
 \qquad |K_R|_{g_R}\leq C/r_y.
\end{equation}
Here and below constants may depend on $L$ and $v$, but not on $R$.

Retain the original forms on the same manifold during the replacement
and bend. Since $y=Ax$ is fixed and all these metrics are uniformly
comparable, their norms are $O(r_y^{-2})$ and their electromagnetic
ball term is $O(r_y^{-4})$. The bend is vacuum for the total
constraints. Thus Lemma~\ref{en:annular-replacement}, together
with $I_{g_R}(Y)=O(r_y^{-4})$, implies the following estimate for
the charged deficit. With $\xi=r_y/R$ there are fixed
$0<a<\xi_0<b$ such that the data are original for $\xi\leq\xi_0$
and horizontal Schwarzschild for $\xi\geq b$, and
\begin{equation}\label{en:charged-deficit}
 \left(-\inf_{|Y|_{g_R}\leq1}D_{g_R}(Y)\right)_+
 \leq
 \begin{cases}
 0,&\xi\leq\xi_0,\\
 \eta_RR^{-3},&\xi_0\leq\xi\leq b,\\
 C r_y^{-4},&\xi\geq b,
 \end{cases}
 \qquad \eta_R\to0.
\end{equation}
The constant radial formulas used next are extended across the core.

\emph{Repair of the charged deficit.}
Choose a smooth $\omega\geq0$ vanishing near $a$ and at least one
on $[\xi_0,b+1]$. On $[a,b+1]$ solve
\[
 z_R'=C_0z_R+\omega,\qquad z_R(a)=0,
\]
and put $z_R=0$ below $a$. The constant $C_0$ will be chosen
large using Equation~\ref{en:radial-bounds}. On the horizontal
tail let
\[
 A_R(t)=(t+m/(2R))^2z_R(t).
\]
Starting at $b$, smoothly continue its derivative to $t^{-2}$
over $[b,b+1]$, using a convex interpolation of the positive
derivative already prescribed and $t^{-2}$. Thus
\begin{equation}\label{en:repair-tail}
 A_R'(t)\geq c t^{-2}\quad(t\geq b),\qquad
 A_R'(t)=t^{-2}\quad\hbox{eventually}.
\end{equation}
The quantities $A_R(\infty)$ and
$F_R(\xi)=\int_\xi^\infty z_R(t)\,dt$ are uniformly bounded
for fixed $L$. Define
\[
 u_R=1+\frac{e_R}{R}F_R(r_y/R),\qquad e_R>0.
\]
It is smooth, at least one, and constant on the compact region
$r_y\leq aR$. Direct differentiation, with $F_R'=-z_R$, gives
\begin{equation}\label{en:repair-radial-calculation}
 \begin{split}
 -\Delta_{g_R}u_R-|K_R||du_R|
 =\frac{e_R}{R^3}\bigl(&z_R'|dr_y|_{g_R}^2
          +Rz_R\Delta_{g_R}r_y\\
          &-R|K_R|z_R|dr_y|_{g_R}\bigr).
 \end{split}
\end{equation}
On $[a,b]$ this is at least
$e_RR^{-3}(c z_R'-C z_R)$. Choose $C_0$ sufficiently large
to make it nonnegative everywhere there and at least
$c'e_RR^{-3}$ on $[\xi_0,b]$.
On the horizontal tail $K_R=0$. Writing
$w_*=1+m/(2r_y)$, the radial conformal Laplacian formula is
$\Delta_{w_*^4\delta}u=w_*^{-6}r_y^{-2}
\partial_{r_y}(r_y^2w_*^2\partial_{r_y}u)$.
It yields the exact identity
\begin{equation}\label{en:repair-schwarzschild}
 -\Delta_{g_R}u_R
 =e_RR^{-3}(1+m/(2R\xi))^{-6}\xi^{-2}A_R'(\xi)
 \geq c''e_RR^{-3}\xi^{-4}.
\end{equation}
Choose $e_R\to0$ with $e_R/(\eta_R+R^{-1})\to\infty$.
For fixed $L$, Equations~\ref{en:charged-deficit},
\ref{en:repair-radial-calculation}, and
\ref{en:repair-schwarzschild}, followed by
Equation~\ref{en:conformal-margin}, show that
$(g_0,K_0)=(u_R^4g_R,u_R^2K_R)$ satisfies charged DEC.
All boundary signs persist, since $u_R$ is constant near $S$.

By Equation~\ref{en:repair-tail},
\[
 u_R=1+\frac{e_RA_R(\infty)}{r_y}+O_\ell(r_y^{-2})
 \quad\hbox{for every }\ell
\]
for each fixed $R,L$. Its ADM energy and momentum are therefore
\begin{equation}\label{en:repaired-mass}
 E_0=m+2e_RA_R(\infty),\qquad P_0=0.
\end{equation}
The second form is compactly supported and the scalar curvature
has $O_\ell(r_y^{-4})$ decay on the tail. Both assertions follow
from the exact horizontal form and
Equation~\ref{en:repair-schwarzschild}.

\emph{Strictification and the limiting comparisons.}
Fix this repaired exterior and apply the compact-$K$ part of
Lemma~\ref{en:strictification}, with a fixed
$0<\delta_0<1$. A further conformal change by $1+s\phi$, $s>0$,
gives the strict inequalities and every stated decay and
integrability property. Its energy tends to $E_0$ as $s\downarrow0$,
its momentum remains zero, and it can be made arbitrarily small
in any prescribed finite collection of compact smooth norms.

It remains to check the area comparison, including cuts that escape
into the bend. On the Schwarzschild graph, the spatial Schwarzschild
metric dominates $|dy|^2$ and its lapse squared is at most one.
Moreover Equation~\ref{en:bend-height} gives
\[
 dT=\psi(v\cdot dy)
        +(v\cdot y)\psi'\frac{dr_y}{Lr_y}.
\]
Since $0\leq\psi\leq1$, the tensor square of this expression is
bounded above by $(v\cdot dy)^2+C L^{-1}|dy|^2$.
Consequently, on the plane and bend,
\begin{equation}\label{en:bend-metric-domination}
 g_R\geq|dy|^2-dT^2
 \geq |dy|^2-(v\cdot dy)^2-C L^{-1}|dy|^2
 \geq(1-C' L^{-1})|dx|^2.
\end{equation}
On the replacement annulus $g_R=\delta_x+O(R^{-\beta})$,
while the original metric has $g=\delta_x+O(R^{-q})$ throughout
the distant end. Inside the replacement radius the metrics agree.
Both conformal changes increase the metric. Therefore the final
metric has a global lower bound
$g_{R,L,s}\geq(1-C'/L-o_R(1))g$.

Choose $L_j\to\infty$ first. For each $L_j$ choose $R_j$ so
large that all the preceding constructions apply, their error
bounds are at most $1/j$, and the energy increment in
Equation~\ref{en:repaired-mass} is at most $1/j$.
The choices are possible because for fixed $L$ the repair constants
and $A_R(\infty)$ are bounded independently of $R$.
Then choose $s_j>0$ small enough that the strictification energy
increment and the errors in the first $j$ compact smooth seminorms
are at most $1/j$. The replacement and bend move to infinity;
$u_R\to1$ uniformly and is constant on every fixed compact set
eventually. Thus $g_j,K_j$ converge locally smoothly to $g,K$.
The field convergence follows from the unchanged forms and the
smooth convergence of the volume forms. This proves the energy,
momentum, and local convergence assertions. The global lower metric
bound proves completeness and gives, for every full cut $\Gamma$,
\[
 \operatorname{Area}_{g_j}(\Gamma)
       \geq(1-\epsilon_j)\operatorname{Area}_g(\Gamma).
\]
Taking infima yields the required lower limit for the areas.
For the upper limit choose, for any $\eta>0$, a fixed smooth full
cut $\Gamma_\eta$ with
$\operatorname{Area}_g(\Gamma_\eta)\leq A_{\min,g}(S)+\eta$.
It is compact, so local convergence gives
\[
 \limsup_j A_{\min,g_j}(S)
 \leq\lim_j\operatorname{Area}_{g_j}(\Gamma_\eta)
 \leq A_{\min,g}(S)+\eta.
\]
Let $\eta\downarrow0$. This proves
Equation~\ref{en:area-convergence} for precisely the original full
cut definition. If one uses the perimeter formulation, Proposition
\ref{ct:perimeter-equivalence} identifies the same infimum, including
the coincident obstacle frontier. At every stage the charge fluxes
were preserved as integrals of the two original closed forms.
\end{proof}
